\sloppy

% Поля по нормоконтролю (верх/низ 20 мм, слева 30 мм, справа 15 мм); класс G7-32 уже подключает geometry — уточняем параметры
\geometry{a4paper,top=20mm,bottom=20mm,left=30mm,right=15mm}

% Абзацный отступ 1,25 см, без дополнительного интервала между абзацами одного стиля (методичка п. 2.2.10)
\setlength{\parindent}{1.25cm}
\setlength{\parskip}{0pt}

% Запрет висячих строк (нормоконтроль п. 35)
\widowpenalty=10000
\clubpenalty=10000

% Настройки стиля ГОСТ 7-32
% Для начала определяем, хотим мы или нет, чтобы рисунки и таблицы нумеровались в пределах раздела, или нам нужна сквозная нумерация.
% \EqInChapter % формулы будут нумероваться в пределах раздела
% \TableInChapter % таблицы будут нумероваться в пределах раздела
% \PicInChapter % рисунки будут нумероваться в пределах раздела

% Добавляем гипертекстовое оглавление в PDF
\usepackage[
bookmarks=true, colorlinks=true, unicode=true,
urlcolor=black,linkcolor=black, anchorcolor=black,
citecolor=black, menucolor=black, filecolor=black,
]{hyperref}

\urlstyle{same}

\usepackage{pdfpages}
\usepackage{pdflscape}

% Изменение начертания шрифта --- после чего выглядит таймсоподобно.
% apt-get install scalable-cyrfonts-tex

% \IfFileExists{cyrtimes.sty}
%     {
        % \usepackage{cyrtimespatched}
%     }
%     {
%         % А если Times нету, то будет CM...
%     }

\usepackage{graphicx}   % Пакет для включения рисунков
\usepackage{grffile}    % поддержка имён файлов с пробелами
\graphicspath{{images/}}

% С такими оно полями оно работает по-умолчанию:
% \RequirePackage[left=20mm,right=10mm,top=20mm,bottom=20mm,headsep=0pt]{geometry}
% Если вас тошнит от поля в 10мм --- увеличивайте до 20-ти, ну и про переплёт не забывайте:

% Пакет Tikz
\usepackage{tikz}
\usetikzlibrary{arrows,positioning,shadows}

% Произвольная нумерация списков.
\usepackage{enumerate}

% ячейки в несколько строчек
\usepackage{multirow}

% itemize внутри tabular
\usepackage{paralist,array}

% Таблицы
\usepackage{amsmath} % equation, \eqref
\usepackage{array} % расширенные возможности для работы с таблицами
\usepackage{tabularx} % автоматический подбор ширины столбцов
\usepackage{dcolumn} % выравнивание чисел по разделителю

\usepackage{setspace} % полуторный интервал (нормоконтроль п. 25)
\setstretch{1.5}

\newenvironment{compactlist}{%
    \renewcommand{\labelenumi}{\hspace{6.5mm}\arabic{enumi})}
    \setlength{\topsep}{0pt}%
    \setlength{\partopsep}{0pt}%
    \setlength{\itemsep}{0pt}%
    \setlength{\parsep}{0pt}%
    \begin{enumerate}%
}{%
    \end{enumerate}%
}



\makeatletter
\renewcommand\theenumi{\arabic{enumi}}
\renewcommand\labelenumi{\hspace{6.5mm}\theenumi)}
\renewcommand\theenumii{\arabic{enumii}}
\renewcommand\labelenumii{\hspace{6.5mm}\theenumii)}
\renewcommand\p@enumii{\theenumi.}
\makeatother

\makeatletter
\renewcommand*\l@section{\@dottedtocline{1}{0em}{2.3em}}
\renewcommand*\l@subsection{\@dottedtocline{2}{0em}{3.2em}}
\renewcommand*\l@subsubsection{\@dottedtocline{3}{0em}{4.1em}}
\makeatother


\usepackage{caption}
\usepackage{ragged2e}

\DeclareCaptionJustification{centerednohyphen}{\Centering\hyphenpenalty=10000\exhyphenpenalty=10000}
\DeclareCaptionLabelSeparator{shortdash}{ \textendash\ }

\usepackage{float}
\setlength{\floatsep}{0pt}
\setlength{\textfloatsep}{0pt}
\setlength{\intextsep}{7pt}

\captionsetup{aboveskip=7pt,belowskip=0pt,justification=centerednohyphen,singlelinecheck=false,labelsep=shortdash}
\captionsetup[table]{justification=RaggedRight,singlelinecheck=false}
\captionsetup[figure]{justification=centerednohyphen,singlelinecheck=false}

\usepackage{chngcntr}

\usepackage{mdframed}
\usepackage{xcolor}

\mdfsetup{
linecolor=black,
linewidth=1pt,
roundcorner=5pt,
innerleftmargin=5pt,
innerrightmargin=5pt,
innertopmargin=5pt,
innerbottommargin=5pt,
skipabove=0pt, % Reduces space above the frame
skipbelow=0pt, % Reduces space below the frame
leftmargin=0pt, % Reduces left outer margin
rightmargin=0pt % Reduces right outer margin
}

% Листинги кода: listings (не minted — не требует pygmentize в PATH при сборке MiKTeX)
\usepackage{listings}
\renewcommand{\lstlistingname}{Листинг}
\lstset{
  basicstyle=\ttfamily\footnotesize,
  breaklines=true,
  columns=fullflexible,
  numbers=left,
  numberstyle=\tiny\color{gray},
  numbersep=8pt,
  stepnumber=1,
  frame=single,
  framerule=0.4pt,
  xleftmargin=2em,
  framexleftmargin=1.2em,
  tabsize=2,
  showstringspaces=false,
  extendedchars=true,
  keepspaces=true
}

\makeatletter
\let\oldtableofcontents\tableofcontents
\renewcommand{\tableofcontents}{%
  \begingroup
  \hyphenpenalty=10000
  \exhyphenpenalty=10000
  \oldtableofcontents
  \endgroup
}
\makeatother

\makeatletter
\def\@hangfrom#1{\setbox\@tempboxa\hbox{{#1}}%
      \hangindent 0pt%\wd\@tempboxa
      \noindent\box\@tempboxa}
\makeatother

\usepackage{etoolbox}
\AtEndEnvironment{lstlisting}{\vspace{-6pt}}
\AtEndEnvironment{table}{\vspace{-10pt}}
\AtEndEnvironment{center}{\vspace{-14pt}}
\AtEndEnvironment{figure}{\vspace{-14pt}}


\usepackage{fontspec}

% Указываем шрифты для основного текста и для кириллицы:
\setmainfont{Times New Roman} % если установлен в системе
\newfontfamily\cyrillicfont{Times New Roman} % для кириллического текста

% Убирает точку после номера главы
\makeatletter
\renewcommand{\thesection}{\thechapter.\arabic{section}}
\renewcommand{\thesubsection}{\thesection.\arabic{subsection}}
\renewcommand{\@seccntformat}[1]{%
  {\fontsize{14pt}{21pt}\selectfont\bfseries \csname the#1\endcsname}\hspace{1em}%
}
\makeatother

\usetikzlibrary{arrows.meta,positioning,shadows,shapes.geometric,calc}
\usetikzlibrary{shapes.multipart}
